\QHeading{第1题　在（　）里填上适当的数。}
\par\noindent


\PlainItem{8.01元 =（{\color{answerred}$80.1$}）角}
\par

\PlainItem{1.5厘米 =（{\color{answerred}$15$}）毫米}
\PlainItem{2.507千米 =（{\color{answerred}$2507$}）米}
\PlainItem{0.35吨 =（{\color{answerred}$350$}）千克}
\PlainItem{9.2千克 =（{\color{answerred}$9200$}）克}
\PlainItem{3.4千克 =（{\color{answerred}$3$}）千克（{\color{answerred}$400$}）克}
\PlainItem{2.04吨 =（{\color{answerred}$2$}）吨（{\color{answerred}$40$}）千克}
\PlainItem{8.76米 =（{\color{answerred}$8$}）米（{\color{answerred}$76$}）厘米}
\PlainItem{6元7角9分 =（{\color{answerred}$6.79$}）元}
\PlainItem{5吨6千克 =（{\color{answerred}$5.006$}）吨}
\par\noindent

\Answer{
\begingroup\color{answerred}
\QHeading{计算过程：1. 单位换算}

\noindent
1. \Formula{\begin{aligned}[t]&\phantom{={}}8.01\text{元}\\&=(8.01\times10)\text{角}\\&=80.1\text{角}\end{aligned}} \par\medskip
 2. \Formula{\begin{aligned}[t]&\phantom{={}}1.5\text{厘米}\\&=(1.5\times10)\text{毫米}\\&=15\text{毫米}\end{aligned}}\par\medskip

3. \Formula{\begin{aligned}[t]&\phantom{={}}2.507\text{千米}\\&=(2.507\times1000)\text{米}\\&=2507\text{米}\end{aligned}} \par\medskip
 4. \Formula{\begin{aligned}[t]&\phantom{={}}0.35\text{吨}\\&=(0.35\times1000)\text{千克}\\&=350\text{千克}\end{aligned}}\par\medskip

5. \Formula{\begin{aligned}[t]&\phantom{={}}9.2\text{千克}\\&=(9.2\times1000)\text{克}\\&=9200\text{克}\end{aligned}} \par\medskip
 6. \Formula{\begin{aligned}[t]&\phantom{={}}3.4\text{千克}\\&=3\text{千克}+0.4\text{千克}\\&=3\text{千克}+(0.4\times1000)\text{克}\\&=3\text{千克}400\text{克}\end{aligned}}\par\medskip

7. \Formula{\begin{aligned}[t]&\phantom{={}}2.04\text{吨}\\&=2\text{吨}+0.04\text{吨}\\&=2\text{吨}+(0.04\times1000)\text{千克}\\&=2\text{吨}40\text{千克}\end{aligned}} \par\medskip
 8. \Formula{\begin{aligned}[t]&\phantom{={}}8.76\text{米}\\&=8\text{米}+0.76\text{米}\\&=8\text{米}+(0.76\times100)\text{厘米}\\&=8\text{米}76\text{厘米}\end{aligned}}\par\medskip

9. \Formula{\begin{aligned}[t]&\phantom{={}}6\text{元}7\text{角}9\text{分}\\&=(6+7\div10+9\div100)\text{元}\\&=(6+0.7+0.09)\text{元}\\&=(6.7+0.09)\text{元}\\&=6.79\text{元}\end{aligned}} \par\medskip
 10. \Formula{\begin{aligned}[t]&\phantom{={}}5\text{吨}6\text{千克}\\&=(5+6\div1000)\text{吨}\\&=(5+0.006)\text{吨}\\&=5.006\text{吨}\end{aligned}}



\endgroup

}

\QHeading{第2题　在$\bigcirc$里填上“$>$”、“$<$”或“$=$”。}

\PlainItem{2.6千克{\color{answerred}$>$}260克}
\par

\PlainItem{3米20厘米{\color{answerred}$=$}3.2米}
\PlainItem{5800平方厘米{\color{answerred}$<$}6平方米}
\PlainItem{4.05吨{\color{answerred}$=$}4吨50千克}
\PlainItem{7元2角{\color{answerred}$>$}7.02元}
\PlainItem{83分米{\color{answerred}$>$}8米3厘米}
\PlainItem{0.75千米{\color{answerred}$<$}760米}
\PlainItem{36厘米{\color{answerred}$<$}3.6米}
\PlainItem{6.56米{\color{answerred}$<$}6米6分米}
\PlainItem{1.2平方米{\color{answerred}$>$}102平方分米}
\par\noindent

\QHeading{第1、2题　单位换算的两种思路}

\Answer{方法一：按进率乘或除；大单位化小单位乘进率，小单位化大单位除进率。方法二：先拆成整数部分和小数部分，分别换算再合并。例如3.4千克可先看成3千克和0.4千克，得到3千克400克。两种方法都必须先确认题目问的是长度、质量还是面积，面积单位进率不能照搬长度单位进率。}

\QHeading{第3题　解决问题。}

\Problem{① 一种矿石每千克3.25元，买100千克和1吨这种矿石各需多少元?}
\par


\MethodTitle{方法一}
\Answer{100千克：$3.25\times100=325$（元）。}
\Answer{1吨的价钱：}
\Answer{\Formula{\begin{aligned}[t]&\phantom{={}}3.25\times(1\times1000)\\&=3.25\times1000\\&=3250\text{（元）}\end{aligned}}}
\Answer{答：购买100千克矿石需要325元。}
\Answer{答：购买1吨矿石需要3250元。}

\MethodTitle{方法二：按每100千克一份计算1吨的费用}
\Answer{100千克是1份，1吨是这样的10份。}
\Answer{\textbf{100千克的价钱}\par \Formula{\begin{aligned}[t]&\phantom{={}}3.25\times 100\\
&=325\text{（元）}\end{aligned}}}
\Answer{\textbf{1吨的价钱}\par \Formula{\begin{aligned}[t]&\phantom{={}}325\times 10\\
&=3250\text{（元）}\end{aligned}}}
\Answer{答：100千克需要325元。\par 答：1吨需要3250元。}
\Problem{② 一个正方形桌面的边长是7分米，现要给桌面的四周镶上铝合金，至少需要铝合金多少米?}

\MethodTitle{方法一}
\Answer{\Formula{\begin{aligned}[t]&\phantom{={}}7\times4\div10\\&=28\div10\\&=2.8\text{（米）}\end{aligned}}}
\Answer{答：至少需要2.8米铝合金。}

\MethodTitle{方法二：先把边长换成米}
\Answer{7分米等于0.7米，再按正方形4条边计算。}
\Answer{\Formula{\begin{aligned}[t]&\phantom{={}}0.7\times 4\\
&=2.8\text{（米）}\end{aligned}}}
\Answer{答：至少需要2.8米铝合金。}


